\documentclass[10pt,a4paper]{amsart}
\usepackage[margin=19mm]{geometry}
\usepackage{amsmath,amssymb,mathtools}
\usepackage[scr=rsfs,cal=euler]{mathalfa}
\usepackage{xfrac}
\usepackage[unicode,hidelinks,bookmarks=false]{hyperref}
\newcommand{\ie}{i.e.\ }
\newcommand{\cf}{cf.\ }
\newcommand{\resp}{respectively\ }
\newcommand{\Subsection}{\subsection}
\providecommand{\nobreakdash}{\nobreak\hbox{-}}
\setlength{\emergencystretch}{2em}
\multlinegap0pt
\usepackage{bm}
\usepackage{bbm}
\usepackage{dsfont}
\usepackage{xspace}
\usepackage{cancel}
\usepackage{mathtools}
\usepackage{amsthm}
\makeatletter
\@ifundefined{theorem}{\newtheorem{theorem}{Theorem}}{}
\@ifundefined{claim}{\newtheorem{claim}{Claim}}{}
\@ifundefined{claimproof}{\newtheorem{claimproof}{Claimproof}}{}
\@ifundefined{lemma}{\newtheorem{lemma}[theorem]{Lemma}}{}
\makeatother
\DeclareMathOperator{\gl}{GL} \DeclareMathOperator{\agl}{AGL}
\def\sub{\subseteq}
\newcommand{\CO}{{\mathcal O}}
\title[Semisimple groups interpretable in various valued fields]{Semisimple groups interpretable in various valued fields}
\author{Yatir Halevi, Assaf Hasson, Ya'acov Peterzil}
\date{Source: arXiv:2309.02727v3 (CC BY 4.0)}
\begin{document}
\maketitle

\noindent\emph{Source and licence.} Semisimple groups interpretable in various valued fields. Version arXiv:2309.02727v3 (CC BY 4.0), distributed under the Creative Commons Attribution 4.0 International licence, as recorded in the arXiv licence metadata for this exact version and shown on the abstract page https://arxiv.org/abs/2309.02727v3. Full rights notice: licenses/arxiv-2309.02727-CC-BY-4.0.txt.\par\smallskip

\noindent\textit{Editorial scope.} Counted unit: the Lemma whose source label is \texttt{end-groups2} in the article (source0.tex lines 768--775, source label \texttt{end-groups2}), delivered together with the article's own complete original proof of it (source0.tex lines 776--826, one \texttt{proof} environment of 51 lines). No mathematics is written, repaired or re-derived: statement and proof are the article's own text line for line. Required context retained uncounted: L:definable subgroup of K or K/O is a ball (lines 693--695); a general lemma (lines 741--744). Every definition, display and label of those lines is retained unchanged. Omitted from this package and not redistributed: 1--692, 696--740, 745--767, 827--2411. Nothing inside a retained excerpt is omitted or altered: every retained body is delivered exactly as the article's lines are written, so no omission receipt of any kind accompanies this package. Citations: the retained excerpts carry no citation command at all, so no bibliography entry of the article is transcribed and no reference is redistributed. Reference adaptation: none was needed and none was made; every reference inside the delivered unit resolves inside it, so no unresolved reference or citation is produced. Printed slips kept literal: no printed word, symbol, hypothesis or number of the counted statement or of its proof was altered. Layout and class: the article's own class is \texttt{amsart} with packages including times, fontenc, amssymb, amsthm, amsmath, mathrsfs, srcltx, todonotes, setspace, geometry, hyperref; the wrapper is the lane's pinned \texttt{amsart} editorial wrapper at 10pt on A4, which restates the theorem-like environments the retained text needs and adds the article's own macro definitions that the retained text needs (\texttt{\textbackslash CO}, \texttt{\textbackslash gl}, \texttt{\textbackslash sub}), with extra wrapper packages (bm, bbm, dsfont, xspace, cancel, mathtools). The delivered numbering is the unit's own. Assets: no retained span contains a figure environment, a graphics-inclusion command, a PDF-page inclusion, a table or a TikZ picture; no image file is copied into this package, the package declares no asset dependency and no image file is redistributed with it. Licence: version arXiv:2309.02727v3 of this article is distributed under the Creative Commons Attribution 4.0 International licence, as recorded in the arXiv licence metadata for this exact version and shown on the abstract page https://arxiv.org/abs/2309.02727v3. A full rights notice is delivered at licenses/arxiv-2309.02727-CC-BY-4.0.txt. Characters: every retained excerpt is pure ASCII, so the delivered engine, XeLaTeX with its default Latin Modern text font, reports no missing character.
\begin{lemma}\label{L:definable subgroup of K or K/O is a ball}
Every  definable subgroup $G$ of $(K,+)$ is a ball, possibly trivial.  As a result, every definable subgroup of $K/\CO$ is a (possibly trivial) ball.
\end{lemma}

\begin{lemma}\label{a general lemma}
Let $I,J, H\subseteq K$ be definable subgroups, $I\subseteq H\cap  J$, and let $T:H/I\to K/J$ be a definable homomorphism.
Then there is $d\in K$ such that, $d\cdot I\subseteq J$ and for every $x\in H$,  $T(x+I)=d\cdot (x+I)+J$.
\end{lemma}

\begin{lemma} \label{end-groups2}
The following holds for all  $n$:

$(1)_n$ If $H\subseteq K^n$ is a definable subgroup of $K^n$ then there is $g\in \gl_n(\CO)$ such that $g(H)$ is a cartesian product  of balls, possibly trivial. 

$(2)_n$ If $H\subseteq K^n$ and $J\subseteq K$ are definable subgroups and $T:H\to K/J$ is a definable homomorphism then there are elements $\alpha_1,\ldots, \alpha_n\in K$ such that for all $x=(x_1,\ldots, x_n)\in H$, 
\[T(x_1,\ldots,x_n)=\alpha_1x_1+\cdots+\alpha_nx_n+J.\]
\end{lemma}

\begin{proof}
$(1)_1$ By Lemma \ref{L:definable subgroup of K or K/O is a ball}, every definable subgroup of $K$ is a ball, possibly trivial.

$(2)_1$ This is  Lemma \ref{a general lemma} for $I=\{0\}$.

We now proceed with the induction step, assuming $(1)_{n-1}, (2)_{n-1}$ and prove $(1)_{n}$: 

Let $\pi:K^n\to K^{n-1}$  be the projection onto the first $n-1$ coordinates. By $(1)_{n-1}$, we may assume that $\pi(H)=H_1\times \cdots \times H_{n-1}$, for balls $H_i\subseteq K$.   Also, write $\ker(\pi)=H\cap (\{0\}^{n-1}\times K)$ as $\{0\}^{n-1}\times J$, for a definable subgroup $J\subseteq K$. 

 Notice that for every $(a,b), (a,c)\in H\sub K^{n-1}\times K$ we have $b-c\in J$ and hence $H$ can be viewed as the graph of a function $T:\pi(H) \to K/J$, mapping $a$ to $b+J$, i.e.
 \[H=\{(a,b)\in K^n:a\in \pi(H)\, \wedge b\in T(a)\}. \]
 By $(2)_{n-1}$, there are $\alpha_1,\ldots, \alpha_{n-1}\in K$, such that $T(x)=\sum_{i=1}^{n-1}\alpha_ix_i+J$.

Hence, 
\[H=\{(x_1,\ldots, x_n)\in K^n: (x_1,\ldots, x_{n-1})\in \pi(H) \wedge x_n-\sum_{i=1}^{n-1} \alpha_i x_i\in J\}. \]
 
 The groups $J$ and $\alpha_iH_i$, for $i=1,\ldots, n-1$, are subgroups of $(K,+)$, hence they are balls. Thus, for every $i=1,\ldots, n-1$,  either $J\sub \alpha_i H_i$ or $\alpha_i H_i\sub J$.  Note that if $\alpha_{i_0} H_{i_0}\sub J$ for some $i_0$  and $(x_1,\dots, x_{n-1})\in \pi(H)$ then $x_n-\sum\limits_{i\neq i_0} \alpha_i x_i\in J$ iff $x_n-\sum\limits_i \alpha_ix_i\in J$.  So there is no harm assuming that $\alpha_i=0$ whenever $J\supseteq \alpha_i H_i$  and that $J\sub \alpha_i H_i$ whenever $\alpha_i\neq 0$. 
Also, we may assume that for some $i$,  $\alpha_i\neq 0$, for otherwise $H=\pi(H)\times J$, and we are done.

Fix $\alpha_1, \dots, \alpha_{n-1}$ as above. Permuting the coordinates, if needed, we may assume that $v(\alpha_1)\leq v(\alpha_j)$, for all $j=2,\ldots, n-1$. Thus, we can write 
\[H=\{(x_1,\ldots, x_n): (x_1,\ldots, x_{n-1})\in \pi(H)\,\wedge \frac{1}{\alpha_1}x_n-(x_1+\sum_{i=2}^{n-1} \frac{\alpha_i}{\alpha_1} x_i)\in \frac{1}{\alpha_1}J\}.\]

Let $S(x_2,\ldots, ,x_n)=\frac{1}{\alpha_1}x_n-\sum_{i=2}^{n-1} \frac{\alpha_i}{\alpha_1}x_i.$ Then $S:K^{n-1}\to K$ is a linear map defined over $\CO$ and we have,

\begin{equation}\label{eq.1}H=\{(x_1,\ldots, x_n): (x_1,\ldots, x_{n-1})\in \pi(H)\,\wedge \, x_1
-S(x_2,\ldots, x_n) \in \frac{1}{\alpha_1}J\}\end{equation}

Let $\hat \pi(x_1,x_2,\ldots, x_n)=(x_2,\ldots, x_n)$ be the projection onto the last $n-1$ coordinates. 
\begin{claim}\label{claim1}
    For every $\hat x=(x_2,\ldots, x_n)\in \hat \pi (H)$, we have $(S(\hat x),\hat x)\in H$.
\end{claim}
\begin{claimproof}
    Let $\hat x=(x_2,\dots,x_n)\in \hat \pi(H)$ and let $x_1=S(\hat x)$, then clearly $x_1-S(\hat x)=0\in \frac{1}{\alpha}J$, so by (\ref{eq.1}), it is sufficient to see that $(x_1,x_2,\ldots, x_{n-1})\in \pi(H)$. Since $\hat x\in \hat \pi(H)$,  there exists $x_1'$ such that $(x_1',x_2,\ldots, x_n)\in H$. In particular, $x_2\in H_2,\ldots, x_{n-1}\in H_{n-1}$, so for $(x_1,\ldots, x_{n-1})$ to be in $\pi(H)$, we only need to verify that $x_1=S(\hat x)\in H_1$. By assumption, $(x_1',x_2,\ldots, x_{n-1}, x_{n})\in H$, so by (\ref{eq.1}), $x_1'\in H_1$ and 
$x_1'-S(\hat x)\in \frac{1}{\alpha_1} J$, so $S(\hat x)\in \frac{1}{\alpha_1}J+x_1'$. However, we assumed that  $J\sub \alpha_1 H_1$ so $\frac{1}{\alpha_1}J\sub H_1$, and therefore $S(\hat x)\in H_1$, hence $(S(\hat x),\hat x)\in H$. 
\end{claimproof}

We get that 
\[H=\{(x_1,x_2,\ldots, x_n):(x_2,\ldots, x_n)\in \hat \pi(H) \,\wedge\, x_1-S(x_2,\ldots, x_n)\in \frac{1}{\alpha_1}J\}.\]


So $H\cap (K\times  \{0\}^{n-1})=\frac{1}{\alpha_1}J\times \{0\}^{n-1}$ and, in particular, the  map $(x_1,\dots, x_n)\mapsto x_1-S(x_2,\dots, x_n)$ from $H$ to $\frac{1}{\alpha_1}J$ is surjective. 
We now define $F:K^n\to K^n$ by \[F(x_1,x_2,\ldots, x_n)=(x_1-S(x_2,\ldots, x_n),x_2, \ldots x_n).\]

Then $F$ is over $\CO$, and by a direct computation one sees that it has determinant $1$, hence $F\in \gl_n(\CO)$. It follows from the definition of $F$ and the observation above that the restriction $F\restriction H$ is definable, injective and onto $\frac{1}{\alpha_1}J\times \hat \pi(H)$. 


By induction, there is $h\in \gl_{n-1}(\CO)$ such that $h(\hat \pi (H))$ is a product of balls. Hence, there is $g\in \gl_n(\CO)$ sending $H$ to a product of balls. This ends the proof of $(1)_n$. 

For $(2)_n$, we start with $T:H\to K/J$. with $H\subseteq K^n$, By $(1)_n$, we may assume that $H=V_1\times \cdots \times V_n$, for definable subgroups $V_i\subseteq K$. Thus, 
\[T(x_1,\ldots, x_n)=T(x_1,0,\ldots, 0)+\cdots +T(0,\ldots, 0,x_n),\] with all elements still in $H$. The result follows from  the case $n=1$.
\end{proof}

\end{document}
